\clearpage
\section{Exhaustive forward operational overlaps and admissible residuals}
\label{sec:copying-forward-overlaps}
This section advances the operational-template argument of
\href{https://arxiv.org/html/1907.01257v5\#A8.SS3.SSS7}{Appendix H.3.7}
towards Proposition~12.2. Every original forward contextual rewrite at a non-entering common-context
position now has a complete local occurrence avoiding shallow micro-beta or
reference-creation holes. This is a classification of actual rewrites, not an assumption about
the shape of a selected match. Arbitrary unrelated multi-edge fillings,
arbitrary sharing trees, recursive boxes and boundary identifications remain
permitted. The original terms, translation and machine rules are unchanged.

A second result supplies the actual pointed residual construction with its
original common-context geometry. Binding freedom is proved for closed sink
replacements and for replacements whose original outputs have binding
ancestors. These results discharge concrete arithmetic, dereference and
atom-equality cases. Full operational robustness still requires the remaining
residual admissibility and local-diagram extraction arguments, the converse
overlap classification, and the universal assembly. Classification of a
source alone does not prove robustness.

\Needspace{11\baselineskip}
\subsection{Why a forward footprint cannot enter a shallow hole}
\termdef{operational-forward-footprint}
Fix an original binding-free common context $C$ and its two recursive filling
families $L,R$. A \emph{forward operational footprint} is a full-label
occurrence $f:A\to C[L]$ of an original local rule source, with its focused
port equal to an original context coordinate. Assume that this rewrite
position is not entering a shallow hole of $C$. The claim is that every edge
of $A$ belongs to a non-hole outer edge of $C$. This concerns only the matched
footprint: unrelated fillings can contain arbitrarily many vertices and edges.

The original left operational graphs have a useful restriction. Their
primitive outer edges consist of the active application or reference head
and stable primitives; a micro-beta source also contains its lambda box.
A stable primitive is passive or an instance edge. In particular, neither
left family has an outer contraction, weakening or store edge. Recursive
bodies inside boxes are unrestricted. These facts survive the independent
recursive isomorphisms in the definition of an operational pair.

At the sole input of each left graph sits its active head. Suppose an expanded
edge consumes a port that came from the original context. Its component
source must be an input: it cannot be a component internal vertex, and graph
well-formedness excludes outputs from source lists. If this component fills
an operational hole, source uniqueness identifies the inspected edge with
the active input head. Thus a passive or boxed head at such a coordinate
cannot come from a shallow operational filling. The proof separately excludes
an unfilled retained hole; it does not assume that every filling certificate
replaces every hole.

The active rule head itself is retained because the focused position is not
entering a hole. Its targets are therefore original context coordinates.
Stable operands attached there avoid holes by the preceding argument.
Sharing and store edges avoid holes by label exclusion. Allocation, read,
assignment, both atom-equality rules, and unary and binary arithmetic have
only these kinds of outer edges. Inspecting their original constructors
therefore gives complete footprint exclusion for all seven families.

For beta, follow the application's function wire to the passive lambda,
then the lambda's target to its actual box. The box label cannot be an active
primitive, so its source at an original coordinate excludes a shallow hole.
The stable argument is covered by the general operand argument. Nothing in
this proof traverses or restricts the lambda body.

For copying, every edge of the original contraction tree is retained by the
label restriction, including siblings off the active resolution path. The
previous binding-path barrier excludes shallow holes throughout the resolved
payload. Its proof may normalize a copyable payload's representation, but
the exclusion is transported back to the supplied original occurrence and
its unchanged tree. No minimal-tree premise is imposed. Finally, extract a
full-label local occurrence from an arbitrary contextual rewrite and apply
the corresponding local classification. This includes original renaming,
plugging and expansion constructors.

\Needspace{11\baselineskip}
\subsection{Preserving the actual common context and resume coordinate}
\termdef{original-edge-residual}
A residual certificate needs more than a relation between its endpoint
graphs. Its common context determines admissibility and its token determines
the continuation. For each untouched original edge, construct a filling
witness whose context is the literal corolla labelled by that edge, including
its actual hole name or complete recursive box body. Retained components
use their existing label-copy isomorphisms. Filled components use the
original operational relation at that very hole. Boxed components retain
the original box and its existing recursive filling data.

Expand these witnesses alongside the local residual in the computed frame.
Their common context is recursively isomorphic to the frame with the local
common context plugged into its distinguished cut edge. In particular,
untouched holes and boxes are not renamed or replaced by arbitrary choices.
This statement concerns the actual original labels; it does not yet prove
the paper's global unique-hole inventory correspondence.

Insert a local common-context port into this assembled context. Commutation
of expansion and interface casts proves that each endpoint position is the
actual corresponding local resume port inserted into the whole successor.
The statement includes internal local ports, not just boundary input zero.
Wire labels are preserved. Consequently binding freedom of the explicit
geometric graft gives a genuine binding-free quasi-specimen at exactly the
two resumed positions. This separates the coordinate theorem, now general,
from admissibility of the graft, which still needs a path argument.

\Needspace{11\baselineskip}
\subsection{Closed sinks and binding-guarded outputs}
\termdef{binding-guarded-outputs}
A \emph{sink} here is a graph with empty output interface. Replacing a complete
sink occurrence by another sink cannot create an adjacency leaving the
inserted component. Every path ending at a surviving hole must therefore
remain among retained edges, where it reflects to the original context.
If the replacement is focus-free and hole-free, no new hole or focus appears
inside it, including recursively. Binding freedom of the original context
then proves binding freedom of the actual graft. The temporary cut context
need not itself be binding-free.

A more general sufficient condition concerns the original source $A$.
It has \emph{binding-guarded outputs} if, for every output, a sourced binding
edge of $A$ is an ancestor of an edge producing that output. There can be no
retained path from such an output to a surviving hole: reflecting that suffix
and prefixing it by the given original binding ancestry would contradict
binding freedom of the original host. For any path in a replacement ending
at a surviving hole, work backwards through its retained suffix. A first
exit from the inserted graph is impossible by the output guard. The whole
path therefore reflects to the original context. This proves admissibility
of any focus-free, hole-free replacement with the same interface. The guard
is a proved sufficient condition, not a new operational rule premise.

Dereference and both atom-equality source graphs satisfy this condition.
Their consumed instance edges are sourced binding edges. The existing
store-resolution traces provide ancestry through each complete atom tree to
the corresponding store edge, and that store produces the relevant output.
Different-store comparison has two outputs and supplies the corresponding
instance path separately for each. The proof does not impose acyclicity on
the surrounding graph or shorten the original atom trees.

\Needspace{11\baselineskip}
\subsection{Actual primitive diamonds and quantitative robustness cases}
\termdef{primitive-residual-diamond}
For a retained footprint containing only primitive labels, the paired right
footprint is recursively isomorphic to the original local source. Indeed,
primitive labels are identical on both sides of each retained certificate;
all carriers and ordered incidences are those of the same matched skeleton.
Collapsing its right corollas supplies this isomorphism and fixes every local
source port. Transport the unchanged local rule through it to construct the
right local rewrite; this step is proved rather than assumed.

Use the reflexive filling witness of the actual local reduct as the local
residual, and insert it into the original common frame. The sink or guarded
output argument proves that this graft is binding-free. Pointed transport
then supplies the residual quasi-specimen. Both sides take one actual rewrite
and resume at the matched successor coordinates. Established machine
successor coherence identifies any independently supplied successor with
this constructed one. Hence the quantitative robustness clause has zero
additional left steps and one right step; reflexivity of the cost order
justifies its cost comparison. Exchanging the same residual gives the
corresponding converse case without assuming symmetry of the operational
relation. These converse cases concern this constructed paired occurrence;
they do not classify every arbitrary converse operational overlap.

Unary negation uses the agreed natural-number meaning $0-n$, and binary
subtraction remains truncated natural subtraction. Neither the arithmetic
semantics nor the original allocation, store, beta or copying rule is
modified by these transfer proofs. General beta/copying residuals can still
contain holes and nontrivial interfaces: the closed replacement arguments
alone do not discharge those remaining cases.

% Declaration correspondence.
\subsection{Lean declarations and proof correspondence}

\Needspace{8\baselineskip}
\paragraph{Left labels and forbidden origins.}
The head discriminator selects application or reference creation. Reading a
stable primitive label preserves its passivity or instance condition by label
injectivity. Induction on the two literal operational sources classifies all
outer primitive labels; the pair theorem transports this through its source
isomorphism. Weakening and contraction contradict both alternatives.
For a general edge certificate, a retained component reflects its label to
the original context, while a filled component has the classified left graph.
This proves forbidden-primitive exclusion, then its sharing and store
specializations. The retained-hole alternative is explicitly ruled out.

\lean{Copying/OperationalLeftLabels}{14}{16}{OperationalTemplates.Kind.head}
\lean{Copying/OperationalLeftLabels}{20}{26}{OperationalTemplates.stable_primitive}
\lean{Copying/OperationalLeftLabels}{31}{64}{OperationalTemplates.Literal.left_primitive_cases}
\lean{Copying/OperationalLeftLabels}{67}{73}{OperationalTemplates.Pair.left_primitive_cases}
\lean{Copying/OperationalLeftLabels}{76}{82}{OperationalTemplates.Pair.left_not_weakening}
\par\penalty0
\lean{Copying/OperationalLeftLabels}{85}{99}{OperationalTemplates.Pair.left_not_contraction}
\lean{Copying/OperationalLeftLabels}{104}{119}{BindingOverlap.forbidden_primitive_not_hole}
\lean{Copying/OperationalLeftLabels}{123}{130}{BindingOverlap.sharing_not_hole}
\lean{Copying/OperationalLeftLabels}{133}{137}{BindingOverlap.store_not_hole}

\Needspace{8\baselineskip}
\paragraph{The active input and actual host coordinates.}
The input-head theorem transports the literal application or reference head
through the pair's source isomorphism. The source-input lemma inspects a
component port: internal ports are disjoint from host ports, and outputs cannot
be sources. The retained-target lemma rewrites the target list as the image of
the original list. At a shallow hole, source uniqueness now forces the active
head label; a retained unfilled hole is excluded by the stated non-hole label
hypothesis. The full-label and primitive corollaries take the contrapositive.

\lean{Copying/OperationalHeadOrigin}{13}{29}{OperationalTemplates.Pair.left_input_head}
\lean{Copying/OperationalHeadOrigin}{32}{39}{Search.Expansion.source_input}
\lean{Copying/OperationalHeadOrigin}{41}{57}{Search.Expansion.Retained.target_origin}
\lean{Copying/OperationalHeadOrigin}{62}{78}{OperationalTemplates.shallow_source_label}
\lean{Copying/OperationalHeadOrigin}{81}{87}{OperationalTemplates.source_not_hole}
\par\penalty0
\lean{Copying/OperationalHeadOrigin}{90}{98}{OperationalTemplates.head_not_hole}

\Needspace{8\baselineskip}
\paragraph{Stable operands.}
Every edge of a stable graph consumes its only input. Map this incidence and
its stable label into the host; a passive or instance label differs from either
active operational head. The host-input exclusion therefore applies to every
edge. The adjacency corollary obtains that input coordinate from a retained
predecessor's target and uses the same argument.

\lean{Copying/OperationalStableOverlap}{14}{29}{OperationalTemplates.ForwardOverlap.stable_avoids}
\lean{Copying/OperationalStableOverlap}{34}{47}{OperationalTemplates.ForwardOverlap.stable_after_retained}

\Needspace{8\baselineskip}
\paragraph{All primitive source constructors.}
The edge-case predicate is the disjunction of being the selected head, an
atom or term sharing edge, a store, or a stable edge adjacent to the head.
Its generic elimination first retains the active head using the non-entering
token, then applies label exclusion or the stable-adjacency lemma to each
remaining alternative. The seven source classifiers inspect the literal
constructor tags. Stable values keep their original stability proofs; each
sharing component uses its arbitrary tree's label theorem. Ordered numeric
arguments and the two store-equality branches are handled separately.

\lean{Copying/OperationalPrimitiveOverlap}{12}{17}{OperationalTemplates.ForwardOverlap.PrimitiveEdgeCase}
\lean{Copying/OperationalPrimitiveOverlap}{22}{44}{OperationalTemplates.ForwardOverlap.primitive_avoids}
\lean{Copying/OperationalPrimitiveOverlap}{48}{65}{OperationalTemplates.ForwardOverlap.allocation_cases}
\lean{Copying/OperationalPrimitiveOverlap}{68}{85}{OperationalTemplates.ForwardOverlap.read_cases}
\lean{Copying/OperationalPrimitiveOverlap}{88}{104}{OperationalTemplates.ForwardOverlap.unary_cases}
\par\penalty0
\lean{Copying/OperationalPrimitiveOverlap}{107}{127}{OperationalTemplates.ForwardOverlap.binary_cases}
\lean{Copying/OperationalPrimitiveOverlap}{131}{159}{OperationalTemplates.ForwardOverlap.assignment_cases}
\lean{Copying/OperationalPrimitiveOverlap}{163}{184}{OperationalTemplates.ForwardOverlap.same_cases}
\lean{Copying/OperationalPrimitiveOverlap}{188}{214}{OperationalTemplates.ForwardOverlap.different_cases}

\Needspace{8\baselineskip}
\paragraph{The complete beta footprint.}
The application head is retained at the original token. Its function target
reaches a stable lambda, whose retained target reaches the actual boxed edge.
Full recursive label agreement excludes every primitive label there, proving
the box's non-hole origin. Every stable argument edge is adjacent to the
application's second target. Splitting the literal outer edge carrier into
head, lambda, box and argument proves the complete claim.

\lean{Copying/OperationalBetaOverlap}{12}{70}{OperationalTemplates.ForwardOverlap.beta_avoids}

\Needspace{8\baselineskip}
\paragraph{The complete copying footprint.}
For a normalized payload, all tree edges avoid holes by the sharing-label
exclusion, and all payload edges do so by the existing binding-path barrier.
Normalize an arbitrary copyable payload only inside the proof, precompose its
occurrence with the inverse representation isomorphism, and transport the
conclusion back using the edge bijection. Other local constructors have heads
incompatible with a contraction. Extracting a local occurrence from a genuine
contextual rewrite therefore gives the copying-specific result.

\lean{Copying/OperationalCopyOverlap}{21}{52}{OperationalTemplates.ForwardOverlap.normal_copy_avoids}
\lean{Copying/OperationalCopyOverlap}{57}{71}{OperationalTemplates.ForwardOverlap.copy_avoids}
\lean{Copying/OperationalCopyOverlap}{75}{91}{OperationalTemplates.ForwardOverlap.local_copy_avoids}
\lean{Copying/OperationalCopyOverlap}{97}{109}{OperationalTemplates.ForwardOverlap.copying_rewrite_occurrence}

\Needspace{8\baselineskip}
\paragraph{Exhaustive original rewrite extraction.}
The local theorem eliminates all nine original constructors, using the seven
primitive classifiers and the separate beta and copying arguments. The
whole-rewrite theorem extracts a local rule with its full-label occurrence
and exact source attachment, then applies that local theorem. The final cut
theorem applies the same proof to the occurrence in the actual successor
factorization. Thus the output includes both shallow-hole avoidance and the
pointed recursive isomorphism to the supplied successor.

\lean{Copying/OperationalForwardOverlap}{21}{48}{OperationalTemplates.ForwardOverlap.local_avoids}
\lean{Copying/OperationalForwardOverlap}{55}{62}{OperationalTemplates.ForwardOverlap.rewrite_occurrence}
\lean{Copying/OperationalForwardOverlap}{67}{73}{OperationalTemplates.ForwardOverlap.rewrite_cut}

\Needspace{8\baselineskip}
\paragraph{Literal corolla witnesses.}
The vertex bijection deletes the impossible internal part of a corolla and
forgets its unique edge index in contributed vertices. The edge bijection
forgets the same index. Their inverse laws use uniqueness of that finite index.
The insertion lemma checks internal, input and output ports; mapping ordered
source and target lists then proves the expansion isomorphism. The original
edge witness uses the literal old label as its context and reuses its retained,
filled or boxed certificate. Its context identity is definitional. Its binding
proof excludes any shallow adjacency in a corolla and inherits recursive
binding freedom from the original edge.

\lean{Copying/OriginalEdgeWitness}{14}{22}{CorollaExpansion.vertices}
\lean{Copying/OriginalEdgeWitness}{24}{31}{CorollaExpansion.edges}
\lean{Copying/OriginalEdgeWitness}{33}{39}{CorollaExpansion.map_insert}
\lean{Copying/OriginalEdgeWitness}{41}{60}{CorollaExpansion.iso}
\lean{Copying/OriginalEdgeWitness}{63}{75}{FillingWitness.originalEdge}
\par\penalty0
\lean{Copying/OriginalEdgeWitness}{77}{79}{FillingWitness.originalEdge_context}
\lean{Copying/OriginalEdgeWitness}{83}{94}{FillingWitness.originalEdge_binding}

\Needspace{8\baselineskip}
\paragraph{Casts and selected-component positions.}
The two cast lemmas eliminate equalities of the endpoint graphs and inspect
the port constructor. The rebounded vertex is the ordinary interface cast
followed by the proved common-context carrier equality. The two position
identities reduce to the cast lemmas, including internal vertices. Position
labels follow from expansion and recursive isomorphism typing, and give the
rebounded label theorem. Finally inspect a selected local port: an internal
port stays inserted, an input becomes the host source port, and an output
becomes the host target port. This proves both endpoint expansion equations.

\lean{Copying/PointedFillingTransport}{10}{19}{FillingWitness.cast_leftPosition}
\lean{Copying/PointedFillingTransport}{21}{30}{FillingWitness.cast_rightPosition}
\lean{Copying/PointedFillingTransport}{32}{35}{FillingWitness.reboundaryVertex}
\lean{Copying/PointedFillingTransport}{37}{47}{FillingWitness.reboundary_leftPosition}
\lean{Copying/PointedFillingTransport}{49}{59}{FillingWitness.reboundary_rightPosition}
\par\penalty0
\lean{Copying/PointedFillingTransport}{61}{65}{FillingWitness.leftPosition_label}
\lean{Copying/PointedFillingTransport}{67}{73}{FillingWitness.reboundaryVertex_label}
\lean{Copying/PointedFillingTransport}{75}{84}{FillingWitness.expand_selected_leftPosition}
\lean{Copying/PointedFillingTransport}{86}{95}{FillingWitness.expand_selected_rightPosition}

\Needspace{8\baselineskip}
\paragraph{Exact geometric graft and pointed residual.}
The component family uses the local reboundary witness at the cut edge and
the literal original-edge witness everywhere else. Its expansion defines the
pointed residual. The geometric graft is the frame with the local common
context plugged in. Extensional equality of its component contexts identifies
the expansion with a single replacement expansion; the existing single-
expansion isomorphism gives the context isomorphism. Binding freedom transports
through this isomorphism. The residual point inserts the rebounded local
common coordinate. Selected-component and reboundary position equations prove
both endpoint identities; expansion typing proves its wire label. Applying the
original binding-free quasi-specimen constructor yields the final theorem.

\lean{Copying/PointedFrameResidual}{17}{21}{RetainedFootprint.originalComponents}
\lean{Copying/PointedFrameResidual}{23}{25}{RetainedFootprint.pointedResidual}
\lean{Copying/PointedFrameResidual}{28}{30}{RetainedFootprint.residualGraft}
\lean{Copying/PointedFrameResidual}{32}{44}{RetainedFootprint.pointedResidual_context}
\lean{Copying/PointedFrameResidual}{48}{51}{RetainedFootprint.pointedResidual_contextIso}
\par\penalty0
\lean{Copying/PointedFrameResidual}{53}{55}{RetainedFootprint.pointedResidual_binding}
\lean{Copying/PointedFrameResidual}{58}{62}{RetainedFootprint.residualPoint}
\lean{Copying/PointedFrameResidual}{64}{72}{RetainedFootprint.residualPoint_left}
\lean{Copying/PointedFrameResidual}{74}{82}{RetainedFootprint.residualPoint_right}
\lean{Copying/PointedFrameResidual}{84}{95}{RetainedFootprint.residualPoint_label}
\par\penalty0
\lean{Copying/PointedFrameResidual}{100}{120}{RetainedFootprint.pointedResidual_binding_quasi}

\Needspace{8\baselineskip}
\paragraph{Sinks and surviving paths.}
The survivor function recovers the original edge from a retained cut edge.
Its label and incidence equations are the defining cut equations; host-host
adjacency reflects through the cut embedding. Empty outputs make any inserted-
to-retained adjacency impossible. Induction backwards along ancestry ending
at a surviving edge shows that the entire path is retained and reflects to
the original graph. A forbidden path must end at a surviving hole because the
replacement is hole-free. Reflect its first sourced binding edge and the
positive path to contradict old binding freedom. The final theorem adds focus
freedom and checks nested boxes by old binding freedom or new hole freedom.

\lean{Copying/BindingSinkReplacement}{18}{19}{CompleteMatchCut.survivor}
\lean{Copying/BindingSinkReplacement}{21}{26}{CompleteMatchCut.survivor_label}
\lean{Copying/BindingSinkReplacement}{28}{34}{CompleteMatchCut.survivor_source}
\lean{Copying/BindingSinkReplacement}{36}{42}{CompleteMatchCut.survivor_target}
\lean{Copying/BindingSinkReplacement}{44}{55}{CompleteMatchCut.survivor_adjacent}
\par\penalty0
\lean{Copying/BindingSinkReplacement}{58}{64}{CompleteMatchCut.sink_no_exit}
\lean{Copying/BindingSinkReplacement}{68}{79}{CompleteMatchCut.sink_ancestor_survivor}
\lean{Copying/BindingSinkReplacement}{83}{106}{CompleteMatchCut.sink_shallow_binding}
\lean{Copying/BindingSinkReplacement}{110}{129}{CompleteMatchCut.bindingFree_replace_sink}

\Needspace{8\baselineskip}
\paragraph{Sink admissibility for actual successors.}
Compare any supplied actual successor with the explicit complete replacement
using established successor coherence, then transport the sink binding theorem
through the resulting graph isomorphism. The frame corollary applies this
geometry to the matched skeleton's occurrence in the original common context.
The unary and binary arithmetic corollaries use their original compute rules
and the focus-free, hole-free numeric reduct.

\lean{Copying/BindingSinkRewrite}{18}{33}{Opening.Occurrence.bindingFree_rewrite_sink}
\lean{Copying/BindingSinkRewrite}{38}{48}{RetainedFootprint.frame_plug_bindingFree_sink}
\lean{Copying/BindingSinkRewrite}{51}{55}{Arithmetic.unary_bindingFree}
\lean{Copying/BindingSinkRewrite}{58}{62}{Arithmetic.binary_bindingFree}

\Needspace{8\baselineskip}
\paragraph{Guards at nonempty output interfaces.}
The guard predicate supplies a sourced binding ancestor for each output.
An inserted-to-retained crossing must use one such original output. Prefixing
a retained suffix to a surviving hole by that original ancestor contradicts
old binding freedom. Backwards induction on a path ending at a surviving hole
therefore rules out its first inserted edge, even with arbitrary boundary
aliases. The shallow-binding theorem reflects the remaining positive path.
The full-binding theorem checks focus freedom and nested boxes exactly as in
the sink proof.

\lean{Copying/BindingGuardedReplacement}{11}{20}{BindingGuardedOutputs}
\lean{Copying/BindingGuardedReplacement}{24}{41}{CompleteMatchCut.guarded_exit_forbidden}
\lean{Copying/BindingGuardedReplacement}{45}{57}{CompleteMatchCut.guarded_ancestor_survivor}
\lean{Copying/BindingGuardedReplacement}{59}{83}{CompleteMatchCut.guarded_shallow_binding}
\lean{Copying/BindingGuardedReplacement}{87}{106}{CompleteMatchCut.bindingFree_replace_guarded}
\par\penalty0

\Needspace{8\baselineskip}
\paragraph{Concrete store guards and actual successors.}
Forget the exact count of a store trace by induction to obtain ancestry.
The original read and equal-store graphs have their consumed instance as a
sourced binding ancestor of the unique store output. Different-store equality
uses its two instance paths for the two outputs. The generic successor theorem
combines guarded replacement with proved successor coherence; the three actual
store-operation corollaries instantiate the original rules and safe reducts.
For primitive retained footprints, the skeleton has identical primitive labels
and identical incidences, so it preserves every output guard. Applying guarded
replacement to that occurrence proves admissibility of the actual frame graft.

\lean{Copying/BindingStoreRewrite}{11}{22}{StoreTrace.ancestor}
\lean{Copying/BindingStoreRewrite}{26}{42}{Dereference.before_bindingGuarded}
\lean{Copying/BindingStoreRewrite}{44}{60}{AtomEquality.Same.before_bindingGuarded}
\lean{Copying/BindingStoreRewrite}{63}{88}{AtomEquality.Different.before_bindingGuarded}
\lean{Copying/BindingStoreRewrite}{91}{97}{Opening.Occurrence.bindingFree_rewrite_guarded}
\par\penalty0
\lean{Copying/BindingStoreRewrite}{100}{105}{Dereference.rewrite_bindingFree}
\lean{Copying/BindingStoreRewrite}{107}{113}{AtomEquality.Same.rewrite_bindingFree}
\lean{Copying/BindingStoreRewrite}{115}{128}{AtomEquality.Different.rewrite_bindingFree}
\lean{Copying/BindingStoreRewrite}{132}{142}{RetainedFootprint.skeleton_bindingGuarded}
\lean{Copying/BindingStoreRewrite}{144}{151}{RetainedFootprint.frame_plug_bindingFree_guarded}
\par\penalty0

\Needspace{8\baselineskip}
\paragraph{A sink residual gives the robustness clause.}
The sink quasi-specimen theorem discharges geometric graft binding freedom
using the closed local residual context. The diamond theorem builds both
original whole-graph rewrites, transports endpoint hole freedom along them,
and attaches that exact residual at the local resume point. The robustness
wrapper then identifies an independently supplied class successor by proved
determinism and supplies zero additional steps against one matching step.

\lean{Copying/SinkFrameResidual}{22}{33}{RetainedFootprint.sink_binding_quasi}
\lean{Copying/SinkFrameResidual}{38}{63}{RetainedFootprint.sink_diamond}
\lean{Copying/SinkFrameResidual}{67}{87}{RetainedFootprint.sink_robust_case}

\Needspace{8\baselineskip}
\paragraph{Primitive reconstruction and sink diamonds.}
The right label-copy is the existing retained corolla on the right. A source
primitive reflects to the original context and then to that right corolla,
proving literal equality of its primitive label. Collapsing these corollas
gives the right reconstruction isomorphism, whose inverse sends each source
port to the corresponding host port of the matched skeleton. Renaming the
original local rewrite through this isomorphism proves the paired local step.
The reflexive witness fixes both endpoint positions. Use it with the sink
diamond to obtain actual primitive whole-graph steps and the quasi-specimen.
The forward and converse quantitative corollaries consume arbitrary supplied
successors; the converse exchanges the same original residual.

\lean{Copying/PrimitiveSinkTransfer}{18}{21}{RetainedFootprint.rightCopy}
\lean{Copying/PrimitiveSinkTransfer}{23}{30}{RetainedFootprint.rightLabel}
\lean{Copying/PrimitiveSinkTransfer}{34}{48}{RetainedFootprint.rightReconstruction}
\lean{Copying/PrimitiveSinkTransfer}{50}{53}{RetainedFootprint.rightReconstruction_inverse}
\lean{Copying/PrimitiveSinkTransfer}{57}{73}{RetainedFootprint.primitive_right_step}
\par\penalty0
\lean{Copying/PrimitiveSinkTransfer}{76}{77}{FillingWitness.refl_leftPosition}
\lean{Copying/PrimitiveSinkTransfer}{79}{94}{FillingWitness.refl_rightPosition}
\lean{Copying/PrimitiveSinkTransfer}{98}{111}{RetainedFootprint.primitive_sink_diamond}
\lean{Copying/PrimitiveSinkTransfer}{115}{130}{RetainedFootprint.primitive_sink_robust_case}
\lean{Copying/PrimitiveSinkTransfer}{134}{150}{RetainedFootprint.primitive_sink_converse_robust_case}
\par\penalty0

\Needspace{8\baselineskip}
\paragraph{Unchanged arithmetic in both directions.}
The unary corollaries instantiate primitive sink transfer with the original
negation source, the numeric result for zero minus the argument, and the
proved primitive footprint and safe reduct. The binary corollaries use the
original operation code, so addition and truncated subtraction share the
same proof. Forward and converse successors are independently quantified.
These generic retained cases explicitly require shallow-hole avoidance;
the operational specializations below prove that premise.

\lean{Copying/ArithmeticResidualRobustness}{19}{32}{Arithmetic.unary_retained_robust}
\lean{Copying/ArithmeticResidualRobustness}{36}{50}{Arithmetic.binary_retained_robust}
\lean{Copying/ArithmeticResidualRobustness}{54}{67}{Arithmetic.unary_retained_converse_robust}
\lean{Copying/ArithmeticResidualRobustness}{69}{83}{Arithmetic.binary_retained_converse_robust}

\Needspace{8\baselineskip}
\paragraph{Guarded primitive diamonds.}
Construct the left whole-graph rewrite by complete-match transfer and the
right rewrite by primitive right reconstruction. The reflexive local reduct
has its actual original common context. Guarded replacement proves binding
freedom of its geometric graft; rewrite preservation gives closed endpoints,
and the pointed residual theorem provides the quasi-specimen. Forward and
converse quantitative cases use the two actual steps and this same residual.
There is no assumed paired rewrite, determinism law or residual certificate.

\lean{Copying/PrimitiveGuardedTransfer}{22}{43}{RetainedFootprint.primitive_guarded_diamond}
\lean{Copying/PrimitiveGuardedTransfer}{45}{60}{RetainedFootprint.primitive_guarded_robust_case}
\lean{Copying/PrimitiveGuardedTransfer}{62}{78}{RetainedFootprint.primitive_guarded_converse_robust_case}

\Needspace{8\baselineskip}
\paragraph{Public operational arithmetic and store cases.}
The five public cases concern unary arithmetic, binary arithmetic, read,
equal-store comparison and distinct-store comparison. Each derives footprint
avoidance from the exhaustive local theorem at the original non-entering token.
The arithmetic cases then apply sink transfer; the store cases apply guarded
primitive transfer with the concrete instance-to-store guards. Output closure
comes from the proved operational pre-template. Consequently callers supply
the actual left occurrence, binding-free common context and closed endpoint
fillings, but no shallow-hole-avoidance, paired-step or residual-admissibility hypothesis.

\lean{Copying/OperationalPrimitiveResiduals}{18}{29}{OperationalTemplates.ForwardResidual.unary}
\lean{Copying/OperationalPrimitiveResiduals}{32}{44}{OperationalTemplates.ForwardResidual.binary}
\lean{Copying/OperationalPrimitiveResiduals}{48}{63}{OperationalTemplates.ForwardResidual.read}
\lean{Copying/OperationalPrimitiveResiduals}{66}{81}{OperationalTemplates.ForwardResidual.same}
\lean{Copying/OperationalPrimitiveResiduals}{85}{102}{OperationalTemplates.ForwardResidual.different}
\par\penalty0
